\subsection{General notions and notations}\label{ssec:notation}

We fix some conventions that will be held all along this paper. If $\mathcal{C}$ is a small category (the class of objects is actually a set) and $\cD$ is any other category, then the symbol $[\mathcal{C}, \cD]$ stands for the category whose objects are functors and whose morphisms are natural transformations between these. Since $\mathcal{C}$ is a small category, the resulting category is in fact a Hom-set category, which means that the class of morphisms (or arrows) between any pair of objects forms a set and this set will be denoted by ${\rm Nat}(F,G)$ for any pair of functors~$F$,~$G$. We shall represent a functor $F \colon \mathcal{C} \to \cD$ between small Hom-set categories as a pair of maps $F=(\Fa,\Fo)$, where $\Fa\colon \Ca \to \Da$ and $\Fo\colon \Co \to \Do$ are the associated maps on the sets of arrows and objects, respectively. Given two objects $d,d' \in \cD$, we denote as usual by $\cD(d,d')$ the set of all arrows from $d$ to $d'$. Assume now that we have a functor $F\colon \mathcal{C} \to \cD$. Then by $\cD(d,\Fa(f))$ we denote the map which sends any arrow $p \in \cD(d, F_{\Sscript{0}}(s(f)))$ to the composition $\Fa(f)p \in \cD(d,F_{\Sscript{0}}(t(f)))$ (here $s(h)$ and $t(h)$ stand for the source and the target of a given arrow $h$). In this way, for each object $d \in \cD$ we have a functor $\cD(d,F(-))$ from $\mathcal{C}$ to the category of sets. Similarly, for each object $d' \in \cD$, we have the functor $\cD(F(-), d')$, as well as the functor $\cD(F(-), F(-))$ from the category $\mathcal{C}^{\Sscript{\rm op}} \times \mathcal{C}$ to the category of sets ($\mathcal{C}^{\Sscript{\rm op}}$ is the opposite category of $\mathcal{C}$ obtained by reversing the arrows of~$\mathcal{C}$).

Let $\Bbbk$ be a fixed base field and $1_{\Sscript{\Bbbk}}$ its identity element. Vector spaces over $\Bbbk$ and their morphisms (i.e.,~$\Bbbk$-linear maps) form a category, denoted by $\Vect$. Finite dimensional ones form a full subcategory of this, denoted by $\vect$. The symbol $\tensor{\Bbbk}$ denotes the tensor product between $\Bbbk$-vector spaces and their $\Bbbk$-linear maps. For any set $S$, we denote by $\Bbbk S:={\rm Span}_{\Sscript{\Bbbk}}\big\{x \,|\, x \in S \big\}$ the $\Bbbk$-vector space whose basis is the set $S$. Any element $x \in S$ is identified with its image $1_{\Sscript{\Bbbk}} x \in \Bbbk S$. By convention $\Bbbk S$ is the zero vector space whenever~$S$ is an empty set. When it is needed, we will also consider~$\Bbbk S$ as a set.

\looseness=-1 In this paper we shall consider rings without identity element (i.e., unity). Nevertheless, we will consider a class of rings (or $\Bbbk$-algebras) which have enough orthogonal idempotents in the sense of \cite{Fuller:1976, Gabriel:1962}, and that are mainly constructed from small categories. Specifically, given any small Hom-set category $\cD$, we can consider the \emph{path algebra} or \emph{Gabriel's ring of $\cD$}: Its underlying $\Bbbk$-vector space is the direct sum $R=\bigoplus_{\Sscript{x, x' \in \Do}} \Bbbk \cD(x,x')$ of $\Bbbk$-vector spaces. The multiplication of this ring is given by
the composition of $\cD$. Thus, for any two homogeneous generic elements $r, r' \in \Da$, the multiplication $(1_{\Sscript{\Bbbk}}r) . (1_{\Sscript{\Bbbk}}r')$ is defined by the rule: $ (1_{\Sscript{\Bbbk}}r) . (1_{\Sscript{\Bbbk}}r') =1_{\Sscript{\Bbbk}} (r r') $, the image of the composition of $r$ and $r'$ when $s(r)=t(r')$, otherwise we set $ (1_{\Sscript{\Bbbk}}r) . (1_{\Sscript{\Bbbk}}r') =0$ (see \cite[p.~346]{Gabriel:1962}). For any $x \in \Do$ we denote by $1_{\Sscript{x}}$ the image of the identity arrow of $x$ in the $\Bbbk$-vector space $R$.

In general, the ring $R$ has no unit, unless the set of objects $\Do$ is finite. Instead of that, it has a set of local units\footnote{Recall that a $\Bbbk$-vector space $R$ endowed with an associative $\Bbbk$-bilinear multiplication is said to be a \emph{ring with local units} over $\Bbbk$ if it has a set of idempotent elements, say $E \subset R$, such that for any finite subset of elements $\{r_1, \dots,r_n\} \subset R$, there is an element $e \in E$ such that $r_ie=er_i=r_i$, for any $i=1,\dots,n$. This means that any two elements $r, r' \in R$ are contained in an unital subring of the form $R_e=eRe$, for some $e \in E$, and $R$ is a~directed limit of the $R_e$'s, see \cite{Abrams:1983, Anh/Marki:1983, Anh/Marki:1987} and~\cite{ElKaoutit:2009}.}. Namely, the local units are given by the set of idempotent elements:
\begin{gather*}
\big\{1_{\Sscript{x_1}} \dotplus \cdots \dotplus 1_{\Sscript{x_n}}\in R \,|\, x_i \in\Do,\, i=1,\dots, n, \text{ and } n \in \mathbb{N}\setminus\{0\}\big\}.
\end{gather*}
For example, if we assume that $\cD$ is a discrete category, that is, the only arrows are the identities (or equivalently $\Do=\Da=X$) then $R=\Bbbk^{(X)}$ is the ring defined as the direct sum of $X$-copies of the base field.

A \emph{unital right $R$-module} is a right $R$-module $M$ such that $MR=M$, \emph{left unital modules} are similarly defined (see \cite[p.~347]{Gabriel:1962}). For instance, the previous ring $R$ attached to the category~$\cD$ decomposes as a direct sum of left and also of right unital $R$-modules:
\begin{gather*}
R = \bigoplus_{x \in \Do} R 1_{\Sscript{x}} = \bigoplus_{x \in \Do} 1_{\Sscript{x}} R.
\end{gather*}
Following \cite{Fuller:1976}, a ring which satisfies these two equalities is referred to as a \emph{ring with enough orthogonal idempotents}, whose complete set of idempotents is the set $\{1_{\Sscript{x}}\}_{x \in \Do}$. A morphism in this category of rings is obviously defined.

The $\Bbbk$-vector space of all homomorphisms between two left $R$-modules $M$ and $M'$ will be denoted by $\hom{R\text{-}}{M}{M'}$. If $T$ is another ring with enough orthogonal idempotents and if $M$ is an $(R, T)$-bimodule ($R$ acts on the left and $T$ on the right), then $\hom{R\text{-}}{M}{M'}$ is considered as left $T$-module by using the standard action $(a .f) \colon M \to M'$ sending $m$ to $f(m a)$ for every $a \in T$ and $f \in \hom{R\text{-}}{M}{M'}$.

\subsection{Description of the main results}\label{ssec:raroes}
A \emph{groupoid} is a small Hom-set category where each morphism is an isomorphism. More precisely, this is a pair of sets
$\cG:=(\cG_{\Sscript{1}}, \cG_{\Sscript{0}})$ with diagram of sets
\begin{gather*}
\xymatrix@C=35pt{
\cG_{\Sscript{1}}\ar@<0.70ex>@{->}|-{\scriptscriptstyle{{s}}}[r]
\ar@<-0.70ex>@{->}|-{\scriptscriptstyle{{t}}}[r] & \ar@{->}|-{
\scriptscriptstyle{\iota}}[l] \cG_{\Sscript{0}}},
\end{gather*}
where as above $s$ and $t$ are the source and the target of a given arrow respectively, and $\iota$ assigns to each object its identity arrow. In addition, there is an associative and unital multiplication $\cG_{\Sscript{2}}:= \cG_{\Sscript{1}} \due \times {\Sscript{{s}}} { \Sscript{{t}}} \cG_{\Sscript{1}} \to \cG_{\Sscript{1}}$ acting by $(f,g) \mapsto fg$, as well as a map $\cG_{\Sscript{1}} \to \cG_{\Sscript{1}}$ which associates to each arrow its inverse. Notice that $\iota$ is an injective map, and so~$\Go$ is identified with a subset of~$\Ga$. Then a groupoid is a category with additional structure, namely the map which sends any arrow to its inverse. We implicitly identify a groupoid with its underlying category. A morphism between two groupoids is just a functor between the underlying categories.

Given a groupoid $\cG$, we denote by $\Rep{\cG}$ its category of $\Bbbk$-linear representations, that is, $\Rep{\cG}=[\cG, \Vect]$, the category of functors from $\cG$ to $\Vect$. Let $\upphi \colon \cH \to \cG$ be a morphism of groupoids and denote by $\upphi_{\Sscript{*}}\colon \Rep{\cG} \to \Rep{\cH}$ the associated restriction functor. The induction and the co-induction functors are denoted by~$\upphi^{*}$ and~${}^{*}\upphi$, respectively (see Lemmas~\ref{lema:Induction} and~\ref{lema:Co-Induction} for the precise definitions of these functors). These are the right and the left adjoint functor of $\upphi_{\Sscript{*}}$. We say that $\upphi$ is a \emph{Frobenius morphism} (see Definition~\ref{def:Frobenius} below) provided that~$\upphi^{*}$ and~${}^{*}\upphi$ are naturally isomorphic.

\looseness=-1 Let $\upphi\colon \cH \to \cG$ be a morphism as above and denote by $A$ and $B$ the rings with enough orthogonal idempotents attached to $\cH$ and $\cG$, respectively. Then there is a $\Bbbk$-linear map given by
\begin{gather}\label{Eq:phi}
\phi\colon \ A \longrightarrow B, \qquad \bigg( \sum_i \lambda_i h_i \longmapsto \sum_i \lambda_i \phia(h_i)\bigg).
\end{gather}

As it is shown in Example~\ref{exam:Phi} below, this map is not in general multiplicative. However, under the assumption that $\phio\colon \Ho \to \Go$ is an injective map, $\phi\colon A \to B$ becomes a homomorphism (or extension) of rings with enough orthogonal idempotents and the complete set of idempotents $\{1_{\Sscript{\phio(u)}}\}_{u \in \Ho}$ is injected into the set $\{1_{\Sscript{x}}\}_{x \in \Go}$. In this way, $B$ becomes an $A$-bimodule via the restriction of scalars, although not necessarily a unital one.

The notion of right (left) groupoid-set and that of groupoid-biset are explicitly recalled in Definitions~\ref{def:Gset} and \ref{def:biset}, respectively. For any morphism $\upphi$ of groupoids as above, we denote $\upg:=\Ga \due \times {\Sscript{s}}{ \Sscript{\upphi_0} } \Ho$ and similarly ${ }^{\Sscript{\upphi}}\uU(\cG):=\Ho \due \times{\Sscript{\upphi_0}}{ \Sscript{t}} \Ga$. These are the right and the left pull-back bisets associated to $\upphi$. More precisely, $\upg$ is the right pull-back $(\cG,\cH)$-biset with structure maps $\varsigma\colon \upg \to \Go$, $(a,u) \mapsto t(a)$ and ${\rm pr}_{\Sscript{2}}\colon \upg \to \Ho$, $(a,u) \mapsto u$. Its right $\cH$-action is given by $(a,u) h= (a\phia(h), s(h))$, whenever $u=t(h)$, and its left $\cG$-action is given by $g(a,u)=(ga,u)$, whenever $s(g)=t(a)$. Similarly, we find that the left pull-back ${ }^{\Sscript{\upphi}}\uU(\cG)$ is an $(\cH,\cG)$-biset with structure maps $\vartheta\colon {}^{\Sscript{\upphi}}\uU(\cG) \to \Go$, $(u,a) \mapsto s(a)$, ${\rm pr}_{\Sscript{1}}\colon \pug \to \Ho, (u,a) \mapsto u$. The left $\cH$-action is given by $h (u,a)=(t(h), \phia(h)a)$, where $s(h)=u$, while the right $\cG$-action is given by $(u,a) g=(u,ag)$, where $s(a)=t(g)$ (see Examples~\ref{exam: HG} and~\ref{exam:bisets} below).

